% Generated draft: requires mathematical and source review.
\documentclass[11pt]{article}
\usepackage[margin=25mm]{geometry}
\usepackage{fontspec}
\setmainfont{Latin Modern Roman}
\usepackage{amsmath,amssymb,mathtools,unicode-math}
\setmathfont{Latin Modern Math}
\usepackage{enumitem,needspace,xurl,hyperref}
\setlength{\parindent}{0pt}
\setlength{\parskip}{5pt}
\setlength{\emergencystretch}{3em}
\newcommand{\N}{\mathbb{N}}
\newcommand{\Z}{\mathbb{Z}}
\newcommand{\Q}{\mathbb{Q}}
\newcommand{\R}{\mathbb{R}}
\newcommand{\C}{\mathbb{C}}
\newcommand{\F}{\mathbb{F}}
\newcommand{\A}{\mathbb{A}}
\newcommand{\PP}{\mathbb{P}}
\newcommand{\E}{\mathbb{E}}
\newcommand{\eps}{\varepsilon}
\newcommand{\dd}{\,\mathrm{d}}
\begin{document}
\section*{UnsolvedMath ID: 701283}
\textbf{Source problem code:} catalogue entry
\subsection{Definitions and mathematical statement}
Let \(X\) be an arbitrary infinite set, let \(F=F(X)\) be the free group with free basis \(X\), and put \(G=\operatorname{Aut}(F)\), the group of all automorphisms of \(F\). The derived subgroup of \(G\) is
\[
G'=[G,G]=\bigl\langle \alpha\beta\alpha^{-1}\beta^{-1}:\alpha,\beta\in G\bigr\rangle.
\]
A group is called perfect if it equals its derived subgroup.

\textbf{Question.} Is \(\operatorname{Aut}(F(X))\) perfect for every infinite set \(X\)? Equivalently, does
\[
\operatorname{Aut}(F(X))=[\operatorname{Aut}(F(X)),\operatorname{Aut}(F(X))]
\]
hold for every infinitely generated free group \(F(X)\)?

The supplied source states that the equality holds when \(X\) is countably infinite, attributing this result to R. Bryant and V. A. Roman'kov, \emph{J. Algebra} \textbf{209} (1998), 713--723. Thus the remaining question concerns arbitrary uncountable ranks. This countable-rank result and the registry's open-status label are supplied source assertions, not independently verified assessments.
\subsection{Short English statement}
Must the automorphism group of every infinitely generated free group be perfect, meaning that every automorphism is a finite product of commutators of automorphisms? The supplied source reports an affirmative answer for countably infinite rank; the question remains directed at uncountable ranks.
\subsection{Sources}
\begin{itemize}
\item Local UnsolvedMath registry, permanent ID \(701283\), exact source code \texttt{catalogue entry}, title \emph{Kourovka Notebook Problem 16.92}. The supplied record labels the problem open and records V. Tolstykh as proposer, with proposed year 2006.
\item \emph{The Kourovka Notebook (Unsolved Problems in Group Theory)}, Problem 16.92, proposed by V. Tolstykh; 16th issue (2006), also identified in the supplied record with the 21st issue edition (2026), arXiv:1401.0300v48. The supplied collection description names E. I. Khukhro and V. D. Mazurov as editors.
\item R. Bryant and V. A. Roman'kov, \emph{J. Algebra} \textbf{209} (1998), 713--723. This reference is supplied for the countably generated case; no article title is supplied, and the reference and its claimed scope have not been independently checked.
\end{itemize}
\end{document}
